\section{อภิปรายและสรุปผล}\label{sec:discussion}
% Thai twin of 05_discussion.tex; keep the two in step. Every number here is reported in Section IV.

บทนี้เชื่อมโยงผลกับงานวิจัยก่อนหน้า ตีความความคลาดเคลื่อนที่ยังเหลืออยู่ เสนอนัยต่อการนำไปใช้
ระบุข้อจำกัดและงานในอนาคต แล้วสรุปผลตามวัตถุประสงค์ทีละข้อ

\subsection{ความเชื่อมโยงกับงานวิจัยก่อนหน้า}\label{sec:prior}
\mbox{DA-v2} Metric-Indoor เรียนหน่วยเมตรจากห้องจำลองของ Hypersim~\cite{da2,hypersim}
และวางห้องทดสอบจริงทุกห้องไกลเกินจริง 1.18 ถึง 1.45 เท่า ZoeDepth~\cite{zoedepth} และ \mbox{DA-v2}
ปรับส่วนหัวแบบเมตริกให้เข้ากับโดเมนเป้าหมาย ด้วยการปรับจูนบนข้อมูลที่มีค่ากำกับ
ผลของเราแสดงว่าวิธีเดียวกันนี้ใช้ได้ เมื่อข้อมูลสอนมาจาก LiDAR ของอุปกรณ์เป้าหมายเอง
เช่นเดียวกับที่ผู้พัฒนา DINOv2 ทำกับงานประมาณความลึก~\cite{dinov2} เราฝึกเฉพาะส่วนถอดรหัสบนส่วนเข้ารหัสที่ตรึงไว้
และรูปทรงของผลทำนายแทบไม่เปลี่ยน (AbsRel แบบ oracle รายเฟรม 0.035 เทียบกับ 0.042)
ซึ่งสอดคล้องกับแนวคิดว่าคุณลักษณะจากส่วนเข้ารหัสถ่ายโอนได้~\cite{yosinski} และส่วนหัวเรียนรู้สเกลเป็นหลัก

ความลึกจากเซนเซอร์ถูกใช้เป็นข้อมูลสอนมานาน ตั้งแต่ Kinect ในอาคารไปจนถึงเครื่องสแกนเลเซอร์บนรถยนต์~\cite{nyuv2,kitti,eigen14}
ความคลาดเคลื่อนราว 1~ซม. ที่ LiDAR ของโทรศัพท์เพิ่มขึ้นจากความลึกของ Faro ในกระบวนการของเรา
สอดคล้องกับความแม่นยำระดับเซนติเมตรที่รายงานไว้เมื่อเทียบกับวิธีทางการสำรวจ~\cite{luetzenburg,spreafico}
และความคลาดเคลื่อนของมันในฐานะข้อมูลสอน (AbsRel \rsLIDARABSREL{}) ต่ำกว่าความคลาดเคลื่อนที่เหลือในโมเดลที่ปรับจูนแล้ว
(\rsFTLIDARALLABSREL{}) ราวหนึ่งอันดับขนาด ซึ่งอาจอธิบายได้ว่าทำไมครูที่เป็นเครื่องสแกนเลเซอร์จึงไม่ได้ให้ผลดีขึ้นอย่างวัดได้
Prompt Depth Anything~\cite{promptda} และงานเพิ่มความละเอียดความลึกของ ARKitScenes~\cite{arkitscenes}
ใช้ LiDAR ขณะใช้งาน ส่วนโมเดลของเราต้องการเพียงภาพสี แต่ต้องแลกด้วยความคลาดเคลื่อนที่สูงกว่ามาก
(\rsFTLIDARALLCM{}~ซม. เทียบกับ \rsLIDARCM{}~ซม.)

โมเดลที่คำนึงถึงกล้อง~\cite{metric3d,unidepth,depthpro} มองว่าความยาวโฟกัสที่ไม่รู้ค่าเป็นอุปสรรคหลักของความลึกแบบเมตริก
แต่สำหรับ Depth Pro บนภาพในอาคารขนาด $640\times480$ นี้ การให้ความยาวโฟกัสจริงไม่ได้ทำให้ใช้งานได้
และความคลาดเคลื่อนหลังตัดสเกลออกยังสูงกว่าของ \mbox{DA-v2} Metric-Indoor หลายเท่า
เราไม่ได้ทดสอบ Metric3D หรือ UniDepth ข้อสรุปนี้จึงไม่ขยายไปถึงแนวทางนี้โดยรวม
เกณฑ์วัดแบบรายภาพ ไม่สามารถแสดงว่าสเกลคงที่ตลอดเฟรมของการสแกนหรือไม่ (บทที่~\ref{sec:related})
และเราพบว่าไม่คงที่ ทั้งในโมเดลความลึกสัมพัทธ์และแบบเมตริก Video Depth Anything~\cite{videoda}
มุ่งแก้ความคงเส้นคงวาข้ามเฟรมของวิดีโอลักษณะนี้ ส่วนจุดเมตริกแบบเบาบาง ที่ใช้ในแนวทางเดียวกับ CNN-SLAM และ MonoSDF~\cite{cnnslam,monosdf}
ให้ผลห่างจาก oracle รายเฟรมเพียง \rsSPARSEGAP{}~ซม. ในรูปแบบที่เป็นขอบบน
ระบบสร้างแบบจำลองที่ใช้การเรียนรู้~\cite{atlas,neuralrecon,simplerecon} ใช้หลายเฟรมต่อการทำนายและใช้ข้อมูลชุดอื่น
ตัวเลขของเราจึงเปรียบเทียบกับของงานเหล่านั้นโดยตรงไม่ได้

\subsection{ทำไมความคลาดเคลื่อนของสเกลยังเหลืออยู่}\label{sec:whyft}
การปรับจูนเรียนรู้ตัวคูณของโดเมนแล้ว \ftmain{} ที่ไม่ได้ปรับเทียบเลย (AbsRel 0.134)
ใกล้เคียงกับโมเดลเดิมที่ได้สเกลและค่าเลื่อนหนึ่งชุดต่อห้องซึ่งหาจากค่าอ้างอิง (0.127)
สิ่งที่การปรับเทียบครั้งเดียวเคยให้ จึงย้ายเข้าไปอยู่ในค่าน้ำหนักของส่วนหัวแล้ว
แต่การแกว่งของสเกลจากเฟรมหนึ่งไปอีกเฟรมไม่ได้ลดลง (1.52 เทียบกับ 1.55)
เหตุผลที่เป็นไปได้มาจากแบบจำลองกล้องรูเข็ม (บทที่~\ref{sec:related}) คือภาพบางแบบ
เช่น ภาพระยะใกล้ของหน้าโต๊ะ หรือภาพผนังเรียบ ดูเหมือนกันที่ระยะต่างกัน และไม่มีสิ่งใดบอกขนาดจริง
โมเดลจึงเดาระยะตามแบบห้องทั่วไป สหสัมพันธ์เชิงลบระหว่างสเกลที่ต้องการกับความลึกของภาพ ซึ่งพบในทุกห้อง สอดคล้องกับคำอธิบายนี้
หากคำอธิบายนี้ถูก การเพิ่มหรือปรับปรุงข้อมูลสอนจะไม่ขจัดการแกว่งนี้ และต้องใช้ข้อมูลเมตริกขณะใช้งาน
จากนั้นการรวมแบบ TSDF~\cite{curless96} จะเฉลี่ยพื้นผิวเดียวกันที่ถูกวางไว้หลายระยะ ให้เป็นพื้นผิวที่หนาและเบลอ
ซึ่งสอดคล้องกับการที่ความคลาดเคลื่อนรายเฟรมระดับปานกลาง กลายเป็น Chamfer distance \rsFTLIDARALLCM{}~ซม.

\subsection{นัยต่อการนำไปใช้}\label{sec:business}
สำหรับแอปบนอุปกรณ์ที่ไม่มี LiDAR ผลชี้ลำดับของงานดังนี้ อันดับแรกคือสเกลเมตริกรายเฟรม
(ห่างจาก oracle รายเฟรมเพียง \rsSPARSEGAP{}~ซม. เทียบกับ \rsSCENECM{}~ซม. เมื่อปรับเทียบครั้งเดียวต่อห้อง
ซึ่งเป็นขอบบน เพราะจุดมาจาก LiDAR ไม่ใช่จาก VIO) ถัดมาคือโมเดลที่ปรับจูนแล้ว ซึ่งขจัดความเอนเอียงของโดเมน
โดยไม่ต้องใช้อะไรเพิ่มขณะใช้งาน และขนาดโมเดลสำคัญน้อยที่สุด (\mbox{DA-v2} Small แย่กว่า Large 1.0~ซม.
เมื่อใช้ oracle รายเฟรม แต่เร็วกว่า 6.4 เท่า) ข้อมูลสอนเก็บได้จากอุปกรณ์รุ่น Pro ระหว่างการใช้งานปกติ
เพราะ ARKit ให้ภาพสี ความลึกจาก LiDAR และท่าทางกล้องมาพร้อมกัน และที่ 3.6 เท่าของความยาววิดีโอบนแล็ปท็อป
กระบวนการในรูปแบบที่ทดสอบจึงทำงานแบบออฟไลน์

\subsection{ข้อจำกัด}\label{sec:limitations}
การประเมินมีขนาดเล็ก คือห้องทดสอบหกห้องในแบบสามมิติ ซึ่งหนึ่งห้องถูกใช้ระหว่างพัฒนา และ 26 ห้องในแบบรายเฟรม
ทั้งหมดเป็นบ้านที่ถ่ายด้วย iPad รุ่นเดียว ไม่มีสำนักงาน ห้องว่าง แสงน้อย หรือกล้องโทรศัพท์รุ่นอื่น
และห้องน้ำที่มีกระจกเงาเป็นห้องที่ยากที่สุดสำหรับทุกแหล่งที่มีสเกลถูกต้อง (ตารางที่~\ref{tab:perroom})
เมื่อมีหกห้อง การทดสอบ Wilcoxon ตรวจพบได้เฉพาะความต่างที่เกิดในทุกห้อง การเปรียบเทียบครูจึงตัดความเป็นไปได้
ของความต่างขนาดเท่าที่เห็นในค่าเฉลี่ย คือหนึ่งถึงสองเซนติเมตร ออกไปไม่ได้
การเปรียบเทียบหลักถูกกำหนดหลังจากประเมินห้องทดสอบทั้งหกแล้ว ห้องอีก \rsVN{} ห้องของการทดลองที่ 7
จึงถูกเพิ่มภายหลัง เพื่อให้ข้อสรุปตั้งอยู่บนห้องที่ไม่ได้ใช้ตัดสินใจเรื่องใดเลยด้วย
แบบจำลองอ้างอิงใช้ท่าทางกล้องจาก ARKit ชุดเดียวกับทุกแถว ระบบจริงบนโทรศัพท์จึงจะมีการเลื่อนของท่าทางกล้องเพิ่มเข้ามา
และความครบถ้วนหมายถึงครบเมื่อเทียบกับส่วนที่ Faro มองเห็นเท่านั้น
งานนี้ปรับจูนเฉพาะส่วนหัวของโมเดลเดียวด้วยสูตรเดียว และ \ftfaro{} กับ \ftlidar{} ต่างกันทั้งจำนวนมุมมองและข้อมูลสอน
checkpoint ถูกเลือกด้วยความลึกจาก Faro รวมถึงของโมเดลที่เรียนจาก LiDAR บนสามห้อง \ftmain{} ถูกเลือกเป็นโมเดลหลัก
เพราะดีที่สุดบนสามห้องนั้น และข้อมูลสอนไม่ต้องใช้เครื่องสแกนเลเซอร์ แต่ลำดับนั้นไม่เกิดซ้ำบนห้องทดสอบ
และ AbsRel ตอนเลือก checkpoint ของโมเดลนี้เปลี่ยนได้ถึง 0.1 ระหว่าง checkpoint ที่ห่างกัน 500 ขั้น
การปรับสเกลครั้งเดียวต่อห้องและด้วยจุดเบาบางเป็นขอบบน เพราะปรับเทียบด้วยค่าอ้างอิง และด้วยจุด LiDAR ที่สะอาดตามลำดับ
Depth Pro รันผ่านโค้ดเพียงชุดเดียว และการตรวจเรื่องความยาวโฟกัสครอบคลุมเพียงห้าเฟรมของห้องเดียว
เกณฑ์ที่ผูก recall กับงานแต่ละแบบเป็นตัวอย่างที่เรากำหนดเอง ไม่ใช่มาตรฐานของอุตสาหกรรม
และเวลาที่วัดได้มาจากโค้ดที่ไม่ได้ปรับแต่ง จึงใช้ได้เฉพาะในรูปอัตราส่วน

\subsection{งานในอนาคต}\label{sec:future}
งานในอนาคตเรียงตามผลกระทบที่คาดไว้ ได้แก่ (1)~ใช้จุดเมตริกรายเฟรมจากระบบติดตามของอุปกรณ์ร่วมกับ \ftmain{}
ซึ่งหัวข้อ~\ref{sec:scale} ให้ AbsRel 0.051 เมื่อทุกเฟรมได้สเกลที่ถูกต้อง
(2)~เลือก checkpoint ด้วยความลึกจาก LiDAR เพื่อไม่ต้องใช้เครื่องสแกนเลเซอร์ในขั้นใดเลย
และเปรียบเทียบครูบนเฟรมชุดเดียวกันกับห้องทดสอบที่มากขึ้น (3)~ฝึกและประเมินกับกล้องโทรศัพท์รุ่นอื่น รวมถึง Android
(4)~ประเมินด้วยจุดเด่นจริงจาก ARKit หรือ ARCore และค่าอ้างอิงที่เก็บด้วยแอปเอง
และ (5)~ทดลองปรับจูนส่วนอื่นของเครือข่ายเพิ่ม รวมถึงโมเดลวิดีโอที่รักษาสเกลให้คงที่ข้ามเฟรม~\cite{videoda}

\subsection{สรุปผล}\label{sec:conclusion}
งานนี้ถามว่าการปรับจูนโมเดลความลึกจากภาพเดี่ยวแบบเมตริก ด้วย LiDAR ของโทรศัพท์เอง
ทำให้แม่นยำกว่าโมเดลเดิมหรือไม่ เมื่อวัดเทียบกับเครื่องสแกนเลเซอร์ สำหรับวัตถุประสงค์ข้อ (1) คำตอบคือใช่
บนห้องที่ทดสอบและตามเกณฑ์ในบทนำ \ftmain{} ลด Chamfer distance จาก \rsMETRICCM{} เหลือ \rsFTLIDARALLCM{}~ซม.
ในห้องทดสอบครบทั้งหก และลด AbsRel จาก \rsVMETRICABSREL{} เหลือ \rsVFTALLABSREL{} ใน \rsVWINS{} จาก \rsVN{} ห้อง
โดยไม่ต้องปรับเทียบหรือใช้เซนเซอร์ขณะใช้งาน และค่าน้ำหนักเรียนจาก LiDAR เพียงอย่างเดียว
แม้ \wone{} จะดีขึ้นเพียงห้าในหกห้องทดสอบ สำหรับข้อ (2) ข้อมูลสอนจากเลเซอร์ไม่ได้ดีกว่าจาก LiDAR อย่างวัดได้
และ 24 ห้องก็ไม่ได้ดีกว่า 10 ห้อง ภายใต้กำลังทางสถิติที่จำกัดของหกห้อง
และการให้ความยาวโฟกัสจริงไม่ได้ทำให้ Depth Pro แข่งขันได้ สำหรับข้อ (3) โมเดลนี้อยู่ระหว่างโมเดลสัมพัทธ์
ที่ปรับเทียบครั้งเดียวต่อห้อง กับที่ปรับด้วยจุดเบาบางทุกเฟรม ไม่รองรับงานใดในสามแบบ
และยังแทน LiDAR ของ iPad (\rsLIDARCM{}~ซม.) ไม่ได้ สำหรับข้อ (4) ความคลาดเคลื่อนที่เหลือคือสเกลที่แกว่งตามมุมมอง
ไม่ใช่การเลื่อนสะสมตามเวลา การปรับจูนขจัดความเอนเอียงของสเกลได้ แต่ไม่ได้ลดการแกว่ง
จุดเมตริกรายเฟรมจากระบบติดตามของอุปกรณ์จึงเป็นขั้นต่อไปที่น่าจะได้ผลที่สุด
แต่การใช้ร่วมกับโมเดลที่ปรับจูนแล้วยังต้องทดสอบต่อไป
